The phase parameters $\delta_x, \delta_y$ take only rational multiples of
$\pi$, tabulated by period in Table~\ref{tab:delta-by-period}. The
zero-phase case $\delta = 0$ is by far the most common, appearing in
$226$ of the $418$ parameter slots. The non-zero phases cluster at
rational multiples of $\pi$ with denominators matching the torus
dimensions: $2\pi/5$ and $4\pi/5$ on the $5 \times 5$, $5 \times 10$, and
$10 \times 10$ tori, and $2\pi/3$ and $4\pi/3$ on the $3 \times 3$ and
$3 \times 6$ tori. This concentration reflects the underlying root-of-unity
structure of the revival condition. Three structural facts follow from
the tables. First, the longest period $N = 60$ occurs on the smallest
prime torus $5 \times 5$. Second, the torus sizes $3 \times 4$, $3 \times
8$, and $4 \times 6$ admit revivals only in the asymmetric configuration
$\rho_x \neq \rho_y$, whereas the remaining tori admit both symmetric and
asymmetric parameter sets. Third, all phase parameters are rational
multiples of $\pi$, with denominators dividing twice the torus dimension.
These three facts together characterise the separable-coin revival
landscape and motivate the algebraic analysis in
Section~\ref{sec:discussion}.
